\documentclass[UTF8]{article}
\usepackage{amssymb,amsmath,amsthm,amscd,latexsym,ctex}
\usepackage{tikz,pgfplots,geometry}
\usepackage{ctex,tikz}
\geometry{a4paper,left=2cm,right=2cm,top=1cm,bottom=1.5cm}
\pgfplotsset{width=10cm,compat=1.16}
\title{大佬称呼不扩散条约}
\date{2020.08.11}
\author{不断电的1528}
\begin{document}
    \maketitle
    % \renewcommand\thesection{\chinese{section}}
    \begin{center}\textbf{
        2020年8月11日在不断电的1528开放供签署\\
        生效日期：2020年8月12日\\
        保存同学：祝至永、章亦流、郑安南和贾海涛}
    \end{center}
    
    \addtocounter{section}{-1}
    \section*{前言}
    本条约各缔约同学（以下称“各缔约同学”），

    考虑到一场大佬战争将使全人类遭受浩劫，因而需要竭尽全力避免发生这种战争的危险并采取措施以保障各同学人民的安全，认为扩散大佬称呼将使发生大佬战争的危险严重增加，根据联合同学大会要求缔结一项防止更广泛地扩散大佬称呼的协定的各项决议，承诺进行合作为应用同学际原子能机构对和平大佬活动的各种保障措施提供便利，表示支持进行研究、发展和其他努力，以促进在同学际原子能机构保障制度范围内应用下述原则，即通过在一定战略地点使用仪器和其他技术有效地保障原料和特殊裂变物质的流动，确认下述原则：
    
    一切缔约同学，不论是有大佬称呼同学家或无大佬称呼同学家，均能为和平目的而获得和平应用大佬技术的利益，包括有大佬称呼同学家由于发展大佬爆炸装置而可能得到的任何技术副产品，深信在促进此项原则时，所有缔约同学均有权参加尽可能充分的科学情报交换，以促进为和平目的而应用原子能的发展，并且单独地或与其他同学家合作，对促进这种发展作出贡献，宣布它们的意图是尽早实现停止大佬军备竞赛和着手采取以大佬裁军为目标的有效措施，敦促所有同学家为达到这个目标而进行合作。

    回顾到1963年禁止在大气层、外层空间和水下进行大佬称呼试验条约的各缔约同学在该条约序言中所表示的谋求:达到永远停止一切大佬称呼,并为此目的继续进行谈判的决心，愿意促进同学际紧张局势的缓和以及各同学间信任的加强，以利于按照在严格和有效的同学际监督下的全面彻底裁军条约停止制造大佬称呼、清除其现有全部储存并从同学家武库中取消大佬称呼及其运载工具。

    回顾到按照联合同学宪章，各同学在其同学际关系中不得使用武力威胁或使用武力，或以不符合联合同学宗旨的任何其他方法侵犯任何同学家的领土完整或政治独立，并且回顾到要尽量减少把世界人力及经济资源用于军备，以促进同学际和平与安全的建立与维护，议定条款如下：
    \newcounter{EX}
    \renewcommand{\theEX}{\text{第}\chinese{EX}\text{条}}
    \newcommand{\Emp}{\vspace{0.3cm}\par\refstepcounter{EX}\textbf{\theEX}\hspace{1em}}
    \Emp 每个有大佬称呼的缔约同学承诺不直接或间接向任何接受同学转让大佬称呼或其他大佬爆炸装置或对这种武器或爆炸装置的控制权；并不以任何方式协助、鼓励或引导任何无大佬称呼同学家制造或以其他方式取得大佬称呼或其他大佬爆炸装置或对这种武器或爆炸装置的控制权。
    \Emp 每个无大佬称呼的缔约同学承诺不直接或间接从任何让与同学接受大佬称呼或其他大佬爆炸装置或对这种武器或爆炸装置的控制权的转让；不制造或以其他方式取得大佬称呼或其他大佬爆炸装置；也不寻求或接受在制造大佬称呼或其他大佬爆炸装置方面的任何协助。
    \Emp \begin{enumerate}
        \item 每个无大佬称呼的缔约同学承诺接受按照同学际原子能机构规约及该机构的保障制度与该机构谈判缔结的协定中所规定的各项保障措施，其目的专为大佬查本同学根据本条约所承担的义务的履行情况，以防止将大佬能从和平用途转用于大佬称呼或其他大佬爆炸装置。原料或特殊裂变物质，无论是正在任何主要大佬设施内生产、处理或使用，或在任何这种设施之外，均应遵从本条所要求的保障措施的程序。本条所要求的各种保障措施应适用于在该同学领土之内、在其管辖之下或在其控制之下的任何地方进行的一切和平大佬活动中的一切原料或特殊裂变物质。
        \item 每个缔约同学承诺不将（a）原料或特殊裂变物质，或（b）特别为处理、使用或生产特殊裂变物质而设计或配备的设备或材料，提供给任何无大佬称呼同学家，以用于和平的目的，除非这种原料或特殊裂变物质受本条所要求的各种保障措施的约束。
        \item 本条所要求的各种保障措施的实施，应符合本条约第四条，并应避免妨碍各缔约同学的经济和技术发展或和平大佬活动领域中的同学际合作，包括按照本条的规定和在本条约序言中阐明的保障原则，为和平目的在同学际上交换大佬材料和处理、使用或生产大佬材料的设备。
        \item 无大佬称呼的缔约同学应单独地或会同其他同学家，按照同学际原子能机构规约与同学际原子能机构订立协定，以适应本条的要求。这类协定的谈判应自本条约最初生效后一百八十天内开始进行。在上述一百八十天期限届满后交存其批准书或加入书的各同学，至迟应自交存之日开始进行这类协定的谈判。这类协定的生效应不迟于谈判开始之日起十八个月。
    \end{enumerate}
    \Emp \begin{enumerate}
        \item 本条约的任何规定不得解释为影响所有缔约同学不受歧视地并按照本条约第一条及第二条的规定开展为和平目的而研究、生产和使用大佬能的不容剥夺的权利。
        \item 所有缔约同学承诺促进并有权参加在最大可能范围内为和平利用大佬能而交换设备、材料和科学技术情报。有条件参加这种交换的各缔约同学还应单独地或会同其他同学家或同学际组织，在进一步发展为和平目的而应用大佬能方面，特别是在无大佬称呼的各缔约同学领土上发展为和平目的应用大佬能方面，进行合作以作出贡献，对于世界上发展中地区的需要应给予应有的考虑。
    \end{enumerate}
    \Emp 每个缔约同学承诺采取适当措施以保证按照本条约，在适当同学际观察下并通过适当同学际程序，使无大佬称呼的缔约同学能在不受歧视的基础上获得对大佬爆炸的任何和平应用的潜在利益，对这些缔约同学在使用爆炸装置方面的收费应尽可能低廉，并免收研究和发展方面的任何费用。无大佬称呼的缔约同学得根据一项或几项特别同学际协定，通过各无大佬称呼同学家具有充分代表权的适当同学际机构，获得这种利益。就此问题的谈判应在条约生效后尽速开始进行。具有这种愿望的无大佬称呼的缔约同学也可以根据双边协定获得这种利益。
    \Emp 每个缔约同学承诺就及早停止大佬军备竞赛和大佬裁军方面的有效措施，以及就一项在严格和有效同学际监督下的全面彻底裁军条约，真诚地进行谈判。
    \Emp 本条约的任何规定均不影响任何同学家集团为了保证其各自领土上完全没有大佬称呼而缔结区域性条约的权利。
    \Emp\begin{enumerate}
        \item 任何缔约同学得对本条约提出修正案。提出的任何修正案应提交各保存同学政府，由各保存同学政府分发给所有缔约同学。随后如经三分之一或三分之一以上缔约同学请求，各保存同学政府应召集会议，邀请所有缔约同学参加，以审议这项修正案。
        \item 本条约的任何修正案须经所有缔约同学的多数票通过，多数票中应包括所有有大佬称呼的缔约同学以及在分发修正案之日系同学际原子能机构理事会理事同学的所有其他缔约同学的票数。该修正案对于每个交存其对该修正案的批准书的缔约同学，应于所有缔约同学的过半数同学家，其中包括所有有大佬称呼同学家以及在分发修正案之日系同学际原子能机构理事会理事同学的所有其他缔约同学，交存其对修正案的批准书时起生效。此后，该修正案对于任何其他缔约同学应在其交存修正案批准书时开始生效。
        \item 本条约生效后五年，应在瑞士日内瓦举行缔约同学会议，审查本条约的实施情况，以保证本条约序言的宗旨和本条约的各项条款正在得到实现。此后，每隔五年，经超过半数缔约同学向各保存同学政府提出以上内容的建议，得另行召集为审查本条约实施情况这一相同目的的会议。
    \end{enumerate}
    \Emp\begin{enumerate}
        \item 本条约应开放供所有同学家签署。凡未在本条约按照本条第三款的规定生效前在本条约上签字的同学家，得随时加入本条约。
        \item 本条约须经签署同学批准。批准书和加入书应交经指定为保存同学政府的大不列颠及北爱尔兰联合王同学、苏维埃社会主义共和同学联盟和美利坚合众同学三同学政府保存。
        \item 本条约应在指定为条约保存同学政府的各同学和本条约的其他四十个签署同学批准本条约并交存其批准书后生效。本条约所称有大佬称呼同学家系指在1967年1月1日前制造并爆炸大佬称呼或其他大佬爆炸装置的同学家。
        \item 对于在本条约生效后交存其批准书或加入书的同学家，本条约应自各该同学交存其批准书或加入书之日起生效。
        \item 各保存同学政府应将每一签字的日期、每份批准书或加入书的交存日期、本条约的生效日期、收到关于召集会议的任何请求的日期以及其他通知事项，迅速告知所有签署同学和加入同学。
        \item 本条约应由各保存同学政府遵照联合同学宪章第一百零二条办理登记。
    \end{enumerate}
    \Emp\begin{enumerate}
        \item 每个缔约同学如果断定与本条约主题有关的非常事件已危及其同学家的最高利益，为行使其同学家主权，应有权退出本条约。该同学应在退出前三个月将此事通知所有其他缔约同学和联合同学安全理事会。这项通知应包括关于该同学认为已危及其最高利益的非常事件的说明。
        \item 本条约生效后二十五年应举行一次会议，以决定本条约是否应无限期地继续有效或应延长一段确定的时期。这种决定应由过半数缔约同学作出。
    \end{enumerate}
    \Emp 本条约的英文、俄文、法文、西班牙文和中文文本具有同等效力；本条约应保存在各保存同学政府的档案库内。各保存同学政府应将经正式大佬证的本条约副本分送各签署同学和加入同学政府。
    
    \begin{center}
        \textbf{下列签署人，经正式授权，在本条约上签字，以资证明.}
    \end{center}
\end{document}